\documentclass[a4paper]{article}

\usepackage{lmodern}
\usepackage[margin=25mm]{geometry}
\usepackage[hidelinks]{hyperref}
\usepackage{array}
\usepackage{xcolor}
\usepackage{verbatim}
\usepackage{corasdiagram}

\title{\texttt{corasdiagram} Manual}
\author{Fabian Robert Smith}
\date{Version 0.1.0}

\begin{document}

\maketitle
\tableofcontents

\section{Overview}

\texttt{corasdiagram} is a notation-first LaTeX package for the five core CORAS
diagram families:

\begin{itemize}
  \item asset diagrams
  \item threat diagrams
  \item risk diagrams
  \item treatment diagrams
  \item treatment overview diagrams
\end{itemize}

The package ships with pre-rendered CORAS icons, semantic diagram
environments, concept-level node macros, and semantic edge macros. The current
supported engines are \texttt{pdflatex} and \texttt{lualatex}.

\section{Installation}

\subsection{Repository-Local Use}

Compile from the repository root with:

\begin{verbatim}
TEXINPUTS=tex/latex//: pdflatex -interaction=nonstopmode -halt-on-error \
  examples/corasdiagram-minimal.tex
\end{verbatim}

\subsection{User TEXMF Tree}

Copy \texttt{tex/latex/corasdiagram/} to
\texttt{\$HOME/texmf/tex/latex/corasdiagram/}. Refresh the TeX filename
database if your TeX distribution requires it.

\subsection{Release Bundle}

The release tooling assembles a CTAN-friendly bundle. Once the package is
published, installation can happen through the normal TeX package manager.

\section{Package Options}

\begin{center}
\begin{tabular}{>{\ttfamily}p{0.22\linewidth}p{0.18\linewidth}p{0.46\linewidth}}
\hline
Option & Default & Meaning \\
\hline
iconset=color|bw & color & Select the color or black-and-white icon set. \\
icons-path=<path> & package icons & Override the icon directory path. \\
\hline
\end{tabular}
\end{center}

\section{Diagram Environments}

\begin{description}
  \item[\texttt{corasassetdiagram}] Stakeholders, assets, asset relations,
    and scope containers.
  \item[\texttt{corasthreatdiagram}] Threat sources, vulnerabilities, threat
    scenarios, unwanted incidents, and impacted assets.
  \item[\texttt{corasriskdiagram}] Threat sources, risks, and impacted assets.
  \item[\texttt{corastreatmentdiagram}] Treatment nodes together with the
    selected threat chain they mitigate.
  \item[\texttt{corastreatmentoverviewdiagram}] Treatments, risks, assets, and
    treatment fan-in via a junction.
\end{description}

Each environment enforces its legal CORAS vocabulary. Illegal node types or
edge relations raise package errors during compilation.

\section{Public API}

\subsection{Canonical Node Macros}

\begin{center}
\begin{tabular}{>{\ttfamily}p{0.34\linewidth}p{0.52\linewidth}}
\hline
Macro & Meaning \\
\hline
\textbackslash corasasset & Asset node with CORAS asset icon. \\
\textbackslash corasstakeholder & Stakeholder node with stakeholder icon. \\
\textbackslash corasthreataccidental & Accidental human threat source. \\
\textbackslash corasthreatdeliberate & Deliberate human threat source. \\
\textbackslash corasthreatnonhuman & Non-human threat source. \\
\textbackslash corasvulnerability & Vulnerability node. \\
\textbackslash corasscenario & Threat scenario node. \\
\textbackslash corasincident & Unwanted incident node. \\
\textbackslash corasunwantedincident & Alias for \textbackslash corasincident. \\
\textbackslash corasrisk & Risk node with title and level metadata. \\
\textbackslash corastreatment & Treatment node. \\
\textbackslash corasjunction & Junction node for treatment overview fan-in. \\
\hline
\end{tabular}
\end{center}

The common keyed API includes \texttt{name}, \texttt{title}, \texttt{meta},
\texttt{level}, \texttt{scope}, \texttt{at}, \texttt{order}, \texttt{layout},
\texttt{right of}, \texttt{below of}, \texttt{incident of}, \texttt{treats},
\texttt{harms}, and \texttt{tikz}.

\subsection{Semantic Edge Macros}

\begin{center}
\begin{tabular}{>{\ttfamily}p{0.28\linewidth}p{0.58\linewidth}}
\hline
Macro & Meaning \\
\hline
\textbackslash corascauses & Solid causal edge with open line arrowhead. \\
\textbackslash corasrelates & Solid relation/impact edge with open line arrowhead. \\
\textbackslash corasconcerns & Dashed concern edge, typically stakeholder to asset. \\
\textbackslash corasassociates & Dashed non-causal association edge. \\
\textbackslash corastreats & Dashed curved treatment edge leaving the source orthogonally. \\
\hline
\end{tabular}
\end{center}

\subsection{Scopes and Containers}

\texttt{\textbackslash corasscope} is the public scope/container helper. The
important keys are:

\begin{itemize}
  \item \texttt{name}
  \item \texttt{scope}
  \item \texttt{fit}
  \item \texttt{kind=asset-scope|system|analysis|risk-group|treatment-scope}
  \item \texttt{padding}
  \item \texttt{stakeholder}
  \item \texttt{stakeholder corner=left|right}
  \item \texttt{tikz}
\end{itemize}

The \texttt{asset-scope} kind renders the rounded asset-diagram container with
the stakeholder corner compartment. The official threat, risk, treatment, and
treatment overview examples do not use containers.

\subsection{Advanced Compatibility API}

The package still ships the lower-level primitives
\texttt{\textbackslash corasnode}, \texttt{\textbackslash corasedge}, and
\texttt{\textbackslash corascontainer}. They remain available for advanced
usage and backward compatibility, but the semantic wrappers are the preferred
public API.

\section{Layout Behavior}

The package supports both explicit coordinates and role-aware automatic
placement. Manual placement always wins when \texttt{at=<...>} is given.

Important layout-related keys:

\begin{itemize}
  \item \texttt{at} for direct placement
  \item \texttt{order} for lane ordering in auto mode
  \item \texttt{scope} for fit-based asset-diagram grouping
  \item \texttt{right of}, \texttt{below of}, \texttt{incident of},
    \texttt{treats}, and \texttt{harms} for relative placement
  \item \texttt{layout=auto|manual}
\end{itemize}

\section{Complete Examples}

\subsection{Asset Diagram}

\begin{center}
\begin{corasassetdiagram}[x=.8cm,y=.8cm]
  \corasstakeholder[
    name=doc-stakeholder,
    scope=doc-asset-scope,
    order=1,
    title={Stakeholder}
  ]
  \corasasset[
    name=doc-asset,
    scope=doc-asset-scope,
    order=1,
    title={Asset}
  ]
  \corasasset[
    name=doc-indirect,
    scope=doc-asset-scope,
    order=2,
    title={Indirect\\Asset}
  ]
  \corasscope[
    name=doc-asset-box,
    scope=doc-asset-scope,
    kind=asset-scope,
    stakeholder=doc-stakeholder,
    stakeholder corner=left
  ]
  \corasrelates[from=doc-asset,to=doc-indirect]
\end{corasassetdiagram}
\end{center}

\subsection{Threat Diagram}

\begin{center}
\begin{corasthreatdiagram}[x=.75cm,y=.75cm]
  \corasthreataccidental[name=doc-employee,order=1,title={Employee}]
  \corasscenario[
    name=doc-remote,
    order=1,
    title={Remote access},
    meta={1 per 6 months}
  ]
  \corasvulnerability[
    name=doc-oldweb,
    order=1,
    title={Old\\webserver}
  ]
  \corasunwantedincident[
    name=doc-disclosure,
    order=1,
    title={Disclosure\\of data}
  ]
  \corasasset[
    name=doc-privacy,
    order=1,
    title={Data\\privacy}
  ]
  \corascauses[from=doc-employee,to=doc-oldweb]
  \corascauses[from=doc-oldweb,to=doc-remote]
  \corascauses[from=doc-remote,to=doc-disclosure]
  \corasrelates[from=doc-disclosure,to=doc-privacy]
\end{corasthreatdiagram}
\end{center}

\subsection{Risk Diagram}

\begin{center}
\begin{corasriskdiagram}[x=.8cm,y=.8cm]
  \corasthreatnonhuman[
    name=doc-infra,
    order=1,
    title={IT-infrastructure}
  ]
  \corasrisk[
    name=doc-risk,
    order=1,
    title={Application\\outage},
    meta={1 per year},
    level=medium
  ]
  \corasasset[
    name=doc-availability,
    order=1,
    title={Application\\availability}
  ]
  \corascauses[from=doc-infra,to=doc-risk]
  \corasrelates[from=doc-risk,to=doc-availability]
\end{corasriskdiagram}
\end{center}

\subsection{Treatment Diagram}

\begin{center}
\begin{corastreatmentdiagram}[x=.75cm,y=.75cm]
  \corasthreataccidental[name=doc-user,order=1,title={Employee}]
  \corasvulnerability[
    name=doc-vulnerability,
    order=1,
    title={Old\\webserver}
  ]
  \corasscenario[
    name=doc-threat,
    order=1,
    title={Remote access},
    meta={1 per 6 months}
  ]
  \corasunwantedincident[
    name=doc-incident,
    order=1,
    title={Disclosure\\of data},
    meta={1 per year}
  ]
  \corasasset[
    name=doc-asset2,
    order=1,
    title={Data\\privacy}
  ]
  \corastreatment[
    name=doc-training,
    order=1,
    title={Security\\training}
  ]
  \corascauses[from=doc-user,to=doc-vulnerability]
  \corascauses[from=doc-vulnerability,to=doc-threat]
  \corascauses[from=doc-threat,to=doc-incident]
  \corasrelates[from=doc-incident,to=doc-asset2]
  \corastreats[from=doc-training,to=doc-threat]
\end{corastreatmentdiagram}
\end{center}

\subsection{Treatment Overview Diagram}

\begin{center}
\begin{corastreatmentoverviewdiagram}[x=.75cm,y=.75cm]
  \corasrisk[
    name=doc-risk1,
    order=1,
    title={Disclosure\\of data},
    meta={direct asset},
    level=unacceptable
  ]
  \corasasset[
    name=doc-asset3,
    order=1,
    title={Data\\privacy}
  ]
  \corasjunction[name=doc-junction]
  \corastreatment[
    name=doc-treatment1,
    order=1,
    title={Limit access}
  ]
  \corastreatment[
    name=doc-treatment2,
    order=2,
    title={Upgrade\\server}
  ]
  \corasrelates[from=doc-risk1,to=doc-asset3]
  \corastreats[from=doc-treatment1,to=doc-junction]
  \corastreats[from=doc-treatment2,to=doc-junction]
  \corastreats[from=doc-junction,to=doc-risk1]
\end{corastreatmentoverviewdiagram}
\end{center}

\section{Source Files}

The repository ships these main author-facing files:

\begin{itemize}
  \item \texttt{tex/latex/corasdiagram/corasdiagram.sty}
  \item \texttt{examples/corasdiagram-minimal.tex}
  \item \texttt{examples/corasdiagram-demo.tex}
  \item \texttt{tests/corasdiagram/*.tex}
\end{itemize}

\section{Build and Release Tooling}

The project includes helper scripts for:

\begin{itemize}
  \item building icon PDFs from the original SVG files
  \item running semantic failure tests
  \item assembling a CTAN-friendly release bundle
  \item building a static documentation site with rendered screenshots
\end{itemize}

See \texttt{CONTRIBUTING.md} for the current repository-root commands.

\end{document}
